% Lösungen Einheit 2 – Nr. 22–32
\einheitenkopf[le2][2 Kreisfläche]{Einheit 2 von 3 · Kreisfläche}
\erg{22}{a) $A = \pi \cdot r^2$ \quad b) $u = \pi \cdot d$ \quad c) $A = \pi \cdot r^2$ \quad d) $u = \pi \cdot d$ \quad e) $A = \pi \cdot r^2$}
\erg{23}{a) 3,1 cm² \quad b) 113,1 cm² \quad c) 153,9 cm² \quad d) 314,2 m² \quad e) $r = 9\,\text{cm}$: 254,5 cm² \quad f) $r = 2{,}7\,\text{cm}$: 22,9 cm²}
\erg{24}{a) 56,5 cm² \quad b) 113,1 cm² \quad c) $r = 7\,\text{cm}$: 77,0 cm² \quad d) $\tfrac{3}{4} \cdot \pi \cdot 4^2 \approx 37{,}7\,\text{cm}^2$}
\erg{25}{a) $\tfrac{1}{4} \cdot \pi \cdot 3^2$ \quad b) $\tfrac{1}{2} \cdot \pi \cdot 4^2$ \quad c) $\tfrac{3}{4} \cdot \pi \cdot 2^2$}
\erg{26}{a) Halbkreis mit $r = 5$ \quad b) Viertelkreis mit $r = 6$}
\erg{27}{Kreis 1: $d = 12\,\text{m}$, $u = 37{,}7\,\text{m}$, $A = 113{,}1\,\text{m}^2$ \quad Kreis 2: $r = 4{,}5\,\text{cm}$, $u = 28{,}3\,\text{cm}$, $A = 63{,}6\,\text{cm}^2$ \quad Kreis 3: $r = 1{,}2\,\text{m}$, $u = 7{,}5\,\text{m}$, $A = 4{,}5\,\text{m}^2$ \quad Kreis 4: $d = 7{,}0\,\text{cm}$, $r = 3{,}5\,\text{cm}$, $A = 38{,}5\,\text{cm}^2$}
\erg{28}{a) $r = \sqrt{154 : \pi} \approx 7{,}0\,\text{cm}$ \quad b) $r = \sqrt{314 : \pi} \approx 10{,}0\,\text{m}$ \quad c) $r \approx 2{,}52\,\text{cm}$, $d \approx 5{,}0\,\text{cm}$ \quad d) $r = 1{,}73\,\text{m}$; $A = \pi \cdot 1{,}73^2 \approx 9{,}40\,\text{m}^2$}
\erg{29}{a) Durchmesser statt Radius eingesetzt; richtig: $r = 6\,\text{cm}$, $A = \pi \cdot 6^2 \approx 113{,}1\,\text{cm}^2$ \quad b) Umfangsformel statt Flächenformel; richtig: $A = \pi \cdot 9^2 \approx 254{,}5\,\text{m}^2$ \quad c) $7^2$ ist $7 \cdot 7 = 49$, nicht $2 \cdot 7$; richtig: $A = \pi \cdot 49 \approx 153{,}9\,\text{cm}^2$}
\erg{30}{a) Nein: Der Radius steht im Quadrat; doppelter Radius heißt $2 \cdot 2 = 4$-fache Fläche. \quad b) Die Hälfte der Stücke liegt oben, die Hälfte unten, jede Hälfte trägt den halben Rand, also ist das Rechteck etwa $\tfrac{u}{2} = \pi \cdot r$ lang; seine Höhe ist die Länge eines Stücks, der Radius. Fläche $\approx \pi \cdot r \cdot r = \pi \cdot r^2$.}
\erg{31}{a) Fläche: $\pi \cdot 0{,}6^2 \approx 1{,}1\,\text{m}^2$ \quad b) Umfang: $\pi \cdot 1{,}2 \approx 3{,}8\,\text{m}$ \quad c) Umfang: $2 \cdot \pi \cdot 3{,}5 \approx 22{,}0\,\text{m}$ \quad d) Fläche: $\pi \cdot 3{,}5^2 \approx 38{,}5\,\text{m}^2$}
\erg{32}{a) groß: $\pi \cdot 15^2 \approx 706{,}9\,\text{cm}^2$; zwei kleine: $2 \cdot \pi \cdot 10^2 \approx 628{,}3\,\text{cm}^2$ \quad b) groß: $8\,€ : 7{,}069 \approx 1{,}13\,€$; klein: $9\,€ : 6{,}283 \approx 1{,}43\,€$ je $100\,\text{cm}^2$ \quad c) Die große Pizza ist günstiger – sie kostet weniger und hat trotzdem mehr Fläche.}
